% Intended LaTeX compiler: pdflatex
\documentclass[a4paper,12pt]{article}
\usepackage[utf8]{inputenc}
\usepackage[T1]{fontenc}
\usepackage{graphicx}
\usepackage{longtable}
\usepackage{wrapfig}
\usepackage{rotating}
\usepackage[normalem]{ulem}
\usepackage{amsmath}
\usepackage{amssymb}
\usepackage{capt-of}
\usepackage{hyperref}
\usepackage{\string~/.emacs.d/common}
\author{mahmood}
\date{\today}
\title{probability homework 3}
\hypersetup{
 pdfauthor={mahmood},
 pdftitle={probability homework 3},
 pdfkeywords={},
 pdfsubject={},
 pdfcreator={Emacs 28.2 (Org mode 9.5.5)}, 
 pdflang={English}}
\begin{document}

\maketitle
homework on \href{20230104221353-random_variable.org}{random variable}s\\
\begin{question}
in an exam there are 10 multiple choice questions and each question has 3 choices of which only one is correct\\
each question amounts to 10 points and the passing score is 60\\
what is the probability for a student who answers randomly to pass the exam\\
\begin{answer}
we find the probabilities for \(6/10,7/10,\dots,10/10\) successes, which are each a \href{20230113174354-binomial_distribution.org}{binomial distribution} and add them\\
\end{answer}
\end{question}
\begin{question}
the probability of a boy being born is 45\%, in a family of 5 kids let \(X\) be the number of boys, what is the distribution of \(X\)?\\
\begin{answer}
\(X\) can equal one of \(0,1,2,3,4,5\), each of which is a binomial distribution\\
\begin{gather*}
p(X=0)=\binom{5}{0}0.45^0{(1-0.45)}^{5-0}=0.0503284\\
p(X=1)=\binom{5}{1}0.45^1{(1-0.45)}^{5-1}=0.205889\\
p(X=2)=\binom{5}{2}0.45^2{(1-0.45)}^{5-2}=0.336909\\
p(X=3)=\binom{5}{3}0.45^3{(1-0.45)}^{5-3}=0.275653\\
p(X=4)=\binom{5}{4}0.45^4{(1-0.45)}^{5-4}=0.112767\\
p(X=5)=\binom{5}{5}0.45^5{(1-0.45)}^{5-5}=0.0184528
\end{gather*}
so the distribution of \(X\) is:\\
\[
p(x) = \begin{cases}
         0.0503284 & X=0\\
         0.205889 & X=1\\
         0.336909 & X=2\\
         0.275653 & X=3\\
         0.112767 & X=4\\
         0.0184528 & X=5
       \end{cases}
\]\\
\end{answer}
\end{question}
\begin{question}
we roll a symmetric die twice\\
X=sum of the results\\
Y=maximum result\\
Z=minimum result\\
find the distribution of each of \(X,Y,Z\)\\
\begin{answer}
\(X\) could be 1-6\\
\(p(X=1)=\frac{1}{36} \quad (1,1)\)\\
\(p(X=2)=\frac{3}{36} \quad (1,2),(2,1),(2,2)\)\\
\(p(X=3)=\frac{5}{36} \quad (1,3),(3,1),(2,3),(3,2),(3,3)\)\\
\(p(X=4)=\frac{7}{36} \quad (1,4),(4,1),(2,4),(4,2),(3,4),(4,4)\)\\
\(p(X=5)=\frac{9}{36} \quad \dots\)\\
\(p(X=6)=\frac{11}{36} \quad \dots\)\\
\(p(X) = \begin{cases} \frac{1}{36} & X=1\\ \frac{3}{36} & X=2\\ \frac{5}{36} & X=3\\ \frac{7}{36} & X=4\\ \frac{9}{36} & X=5\\ \frac{11}{36} & X=6 \end{cases}\)\\
\(Y\) could be 1-6\\
\(p(Y=1)=\frac{11}{36} \quad (1,1),(1,2),(2,1),(1,3),(3,1),(1,4),(4,1),(1,5),(5,1),(1,6),(6,1)\)\\
\(p(Y=2)=\frac{9}{36} \quad (2,2),(2,3),(3,2),(2,4),(4,2),(2,5),(5,2),(2,6),(6,2)\)\\
\(p(Y=3)=\frac{7}{36} \quad \dots\)\\
\(p(Y=4)=\frac{5}{36} \quad \dots\)\\
\(p(Y=5)=\frac{3}{36} \quad \dots\)\\
\(p(Y=6)=\frac{1}{36} \quad \dots\)\\
\(p(Y) = \begin{cases} \frac{11}{36} & Y=1\\ \frac{9}{36} & Y=2\\ \frac{7}{36} & Y=3\\ \frac{5}{36} & Y=4\\ \frac{3}{36} & Y=5\\ \frac{1}{36} & Y=6 \end{cases}\)\\
\(Z\) could be 2-12\\
\(p(Z=2)=\frac{1}{36} \quad (1,1)\)\\
\(p(Z=3)=\frac{2}{36} \quad (1,2),(2,1)\)\\
\(p(Z=4)=\frac{3}{36} \quad (2,2),(3,1),(1,3)\)\\
\(p(Z=5)=\frac{4}{36} \quad (2,3),(3,2),(1,4),(4,1)\)\\
\(p(Z=6)=\frac{5}{36} \quad (3,3),(4,2),(2,4),(5,1),(1,5)\)\\
\(p(Z=7)=\frac{6}{36} \quad (3,4),(4,3),(5,2),(2,5),(6,1),(1,6)\)\\
\(p(Z=8)=\frac{5}{36} \quad (4,4),(3,5),(5,3),(6,2),(2,6)\)\\
\(p(Z=9)=\frac{4}{36} \quad \dots\)\\
\(p(Z=10)=\frac{3}{36} \quad \dots\)\\
\(p(Z=11)=\frac{2}{36} \quad (5,6),(6,5)\)\\
\(p(Z=12)=\frac{1}{36} \quad (6,6)\)\\
\begin{gather*}
p(Z) = \begin{cases}
            \frac{1}{36} & Z=2\\
            \frac{2}{36} & Z=3\\
            \frac{3}{36} & Z=4\\
            \frac{4}{36} & Z=5\\
            \frac{5}{36} & Z=6\\
            \frac{6}{36} & Z=7\\
            \frac{5}{36} & Z=8\\
            \frac{4}{36} & Z=9\\
            \frac{3}{36} & Z=10\\
            \frac{2}{36} & Z=11\\
            \frac{1}{36} & Z=12
          \end{cases}
\end{gather*}
\end{answer}
\end{question}
\begin{question}
a pitcher contains 8 balls, 3 of which are red, we withdraw balls one at a time without returning, what is the distribution of the \href{20230104221353-random_variable.org}{random variable} \(X\) that denotes the number of reds that have to be removed until a non-red ball is removed\\
\begin{answer}
i have no idea what is meant here\\
\end{answer}
\end{question}
\begin{question}
on the sides of a coin written are the numbers -1, +1, we throw the coin \texttt{n} times, let \(Y\) be the sum of the numbers we got, if the probability of a +1 equals p, what is the distribution of \(Y\)?\\
\begin{answer}
\(Y\) would vary from \(-n\) to \(+n\)\\
probability of the sum being 2k-n = probability of k +1's\\
\begin{note}
didnt solve this one aswell :(\\
\end{note}
\end{answer}
\end{question}
\end{document}